\section{Boundaries}

Spatial outliers reflect real geography as often as they reflect data problems.
Neighboring units can differ because they sit on opposite sides of a city line,
county border, campus boundary, reservation boundary, commuting corridor, or
housing-market break. A W.05 score is therefore descriptive. It says that a unit
is unusual relative to its local spatial baseline; it does not say why.

The neighbor definition drives the result. Shared-boundary adjacency, rook
versus queen-style contact rules, minimum boundary length, islands, water
crossings, enclaves, precinct splits, and missing geometries can all change the
neighbor set. Edge units also have fewer neighbors and can receive unstable
scores. The report must carry enough adjacency metadata and investigation notes
for a reviewer to see those choices.

The claim boundary is the same one stated at the beginning. Spatial forensic
analytics can prioritize investigation. They cannot prove fraud, certify an
election, or replace a risk-limiting audit. Any public use of W.05 should keep
its transcript separate from V-series certification evidence and should present
high scores as questions for review, not answers.
